\section{Conditional kernels and ordered paths}\label{sec:paths}

Fix the good stage supplied by Equation~\ref{eq:good-stage-signal}.
In this section abbreviate $s=s_j$, $\alpha=\alpha_j$, $\beta=\beta_j$,
$I=I_j$, $S=S_j$, and $J=J_{j-1}$. All background experiments are
through stage $j$ only. We first identify the conditional laws needed
to turn amplitude variation into a common response on predicates.

\subsection{The deleted-slot background}

Fix an exterior value $\cZ=z$. It fixes the freeze sets, the complete
frozen variables, and the sets $I,S$. Let $Y_i=(\cR_i,a_i)$ for $i\in I$,
with common law $P_Y=P_{\cR}\otimes\lambda$. Let $X$ contain all
unfrozen primary samples and actual trial variables at stages before
$j$, all future primary samples, and the fixed frozen contributions.
It determines the preceding partition sequence
$C^-=C_{\le j-1}$ and the score $B_{\rm out}$ contributed by slots
outside $I$ after earlier stages. It includes no analytic reference
gradients or signs as inputs. Actual previous trials used only their
preceding query partitions to interpret independent descriptors, so
$X$ is independent of the current unfrozen primary samples conditional
on $z$.

For a preceding partition sequence $C$, write the current trial kernel as
\begin{equation}\label{eq:trial-kernel}
 q_C=(1-\alpha)\delta_{\rm shared}
 +\frac\alpha L\sum_{m=1}^L
       P_m(\dd E)\operatorname{Unif}(\{-1,+1\}^{J})(\dd\tau)
          \operatorname{Unif}[0,1](\dd u)
          \operatorname{Unif}\{-1,+1\}(\dd\chi),
\end{equation}
where $P_m$ is the fresh data law of mode $m$, including its fresh
amplitude if required. The descriptor in this formula is interpreted
on $C_{j-1}$. If $T_i$ denotes the current trial, the exact conditional
factorization is
\begin{equation}\label{eq:conditional-product}
 P_z(\dd X,\dd Y,\dd T)
 =q_z(\dd X)\prod_{i\in I}
                  [P_Y(\dd Y_i)q_{C^-(X)}(\dd T_i)].
\end{equation}
This conditional-on-$X$ product is the independence statement we use.
After mixing over $X$, no mutual independence of interpreted trials
is needed.

Let $V_C(Y_i,T_i)$ be the current slot summand, including its factor $s$.
Then
\[
 B^j_{\neg i}=B_{\rm out}+\sum_{k\in I\setminus\{i\}}V_{C^-}(Y_k,T_k),
 \qquad
 B_\circ^j=B_{\rm out}+\sum_{k\in I\setminus S}V_{C^-}(Y_k,T_k).
\]
For every $i\in I$, the pair $(B^j_{\neg i},C^-)$ uses none of $Y_i$.
Equation~\ref{eq:conditional-product} shows that its conditional law
given $z$ is independent of the entire primary presentation and amplitude
in that slot. This law is also the same for every $i\in I$: permutations
of these slots preserve the product law and send one deleted sum to
another. Denote the common law by $\kappa_z(\dd B,\dd C)$.
The actual fine partition $C_j$ is not a coordinate of this kernel.

Write $\phi_C(B)$ for lookup smoothed through the preceding partition
sequence $C$. In the experiment forcing $i$ shared, put
$F_i=\phi_{C^-}(B^j_{\neg i}+sa_ir_i)$, where $r_i$ is the predicate
represented by $\cR_i$. The definition of $b_i$ and the tower property give
\[
 s g_i=\E^{i,\mathrm{sh}}[F_i\mid\cZ,\cR_i,a_i]
              -\E^{i,\mathrm{sh}}[F_i\mid\cZ,\cR_i].
\]
Indeed $\sigma(\cZ,\cR_i,a_i)\subseteq\cB\vee\sigma(a_i)$, and
$a_i$ is independent of $\cB$; averaging the subtracted
$\cB$-conditional expectation therefore removes $a_i$. Forcing shared
preserves the law of the background. The projected response is consequently
\begin{equation}\label{eq:common-amplitude-kernel}
 s\,g_z(r,a)=\int\left[
       \phi_C(B+sar)-\int\phi_C(B+sa'r)\dd\lambda(a')
                      \right]\dd\kappa_z(B,C).
\end{equation}
It is a version of every $g_i$, evaluated at the predicate represented
by $\cR_i$ and at $a_i$. It depends on the actual predicate $r$, not
its presentation. In particular $g_z(0,a)=0$. An unused queried position
in a zero-predicate presentation can change $C_j$, but it cannot change
Equation~\ref{eq:common-amplitude-kernel}, which uses only the preceding
partitions. Independent sampled predicates may have overlapping bit
names; independence here concerns samples and does not require disjoint
supports on the assignment cube.

For $i\in S$, extend this background by
$\sigma=\operatorname{sign}K_{j-1}(B_\circ^j)$.
The triple $(B^j_{\neg i},C^-,\sigma)$ still uses none of $Y_i$ and is
unaffected by changing $T_i$. Its conditional law is common to all
$i\in S$, by permutations within $S$; call it $\kappa_z^S$.
Its $(B,C)$ marginal is precisely $\kappa_z$.
This extended independence is asserted only for $i\in S$.
The sign can be correlated with the remaining background, and indices
outside $S$ can enter its definition.

Define the coarse width on any actual predicate $r$ by
\begin{equation}\label{eq:coarse-width}
 F_z(r)=\int\frac{\phi_C(B+s\sigma r)-\phi_C(B-s\sigma r)}s
                       \dd\kappa_z^S(B,C,\sigma).
\end{equation}
Fresh $r$ is independent of this background. Averaging over uniform
$i\in S$ in the ordinary experiment is another description of this
integral. Because its two amplitude terms ignore $\sigma$, the same
extended kernel also represents Equation~\ref{eq:common-amplitude-kernel}.
This supplies one common background for all three values at
$ar,-\sigma r,+\sigma r$.

Figure~\ref{fig:conditional-kernels} records the two conditional laws
and their common marginal.

\begin{figure}[tbp]
\centering
\begin{tikzpicture}[
  font=\small,
  kernel/.style={draw=black!55,fill=blue!3,rounded corners=2pt,
    align=center,text width=6.1cm,minimum height=1.35cm,inner sep=5pt},
  response/.style={draw=black!40,rounded corners=2pt,
    align=center,text width=6.1cm,minimum height=1.2cm,inner sep=5pt},
  map/.style={-{Stealth[length=2mm]},draw=black!70,thick}
]
\node[align=center] at (0,1.2)
  {Fix the exterior $\cZ=z$; write $Y_i=(\cR_i,a_i)$.};
\node[kernel] (ordinary) at (-3.95,0) {
  \textbf{Every $i\in I$}\\[3pt]
  $(B^j_{\neg i},C^-)\mid(z,Y_i)\sim\kappa_z$\\[3pt]
  The current fine partition $C_j$ is absent.
};
\node[kernel] (extended) at (3.95,0) {
  \textbf{Only $i\in S$}\\[3pt]
  $(B^j_{\neg i},C^-,\sigma)\mid(z,Y_i)\sim\kappa_z^S$\\[3pt]
  $\sigma=\operatorname{sign}K_{j-1}(B_\circ^j)$
};
\draw[map] (extended.west) -- node[above,font=\scriptsize]
  {forget $\sigma$} (ordinary.east);
\node[response] (amplitude) at (-3.95,-2.0) {
  $s\,g_z(r,a)=\E_{\kappa_z}\big[\phi_C(B+sar)$\\[3pt]
  $-\E_{a'\sim\lambda}\phi_C(B+sa'r)\big]$
};
\node[response] (width) at (3.95,-2.0) {
  $s\,F_z(r)=\E_{\kappa_z^S}\big[\phi_C(B+s\sigma r)$\\[3pt]
  $-\phi_C(B-s\sigma r)\big]$
};
\draw[map] (ordinary.south) -- node[right,font=\scriptsize]
  {amplitude response} (amplitude.north);
\draw[map] (extended.south) -- node[right,font=\scriptsize]
  {signed width} (width.north);
\end{tikzpicture}
\caption{The common coarse backgrounds, conditional on $\cZ=z$.
The displayed laws do not depend on $Y_i$; their index ranges differ.
Deleting the whole subbatch $S$ in $B_\circ^j$ makes the reference
sign independent of each tested $Y_i$, conditional on $z$, while
allowing dependence on the remaining background. The $(B,C)$ marginal
of $\kappa_z^S$ is $\kappa_z$, so the amplitude response can also be
averaged over the extended kernel. Here $C=C^-$, $\phi_C$ uses only
preceding smoothing, and the fresh predicate and amplitudes are
independent of the background conditional on $z$.}
\label{fig:conditional-kernels}
\end{figure}

\subsection{Forcing a slot and rebuilding the fine state}

Fix a mode $m$ and an index $i\in S$. Let $\cH_i$ contain $z$, $X$,
every current primary pair $Y_k$, and all current trials $T_k$ except
$T_i$. In particular it retains $Y_i$, whose original question list
still enters the fine partition when its summand is replaced.
The preceding partitions, $B_\circ^j$, and $\sigma$ are
$\cH_i$-measurable. The event
\[
 E_i=\{T_i\text{ is replacement in mode }m,
                  \ \tau_i=\sigma,\ \chi_i=+1\}
\]
has the exact conditional probability
\begin{equation}\label{eq:match-probability}
 \Prob(E_i\mid\cH_i)=\frac\alpha{2L2^J}=\beta.
\end{equation}
The number of formal parent cells is deterministic; empty cells and
colliding labels have not reduced $J$. Conditional on $E_i,\cH_i$, the
path data have their original mode law and $u$ is independent uniform.

Let $\Gamma(\cH_i,T_i)$ construct the entire stage-$j$ state, including
the score and the fine partition rebuilt from all unchanged primary
question lists and the trial's new path question list. For every
nonnegative measurable full-state function $Q_0$, the forced experiment
preserves the ordinary marginal law of $\cH_i$ and satisfies
\begin{equation}\label{eq:full-state-forcing}
 \E_{\rm forced}[Q_0(\Gamma)\mid\cH_i]
 =\E_{\rm actual}
       [Q_0(\Gamma)\ind_{E_i}/\beta\mid\cH_i].
\end{equation}
This is just conditional expectation under the event in
Equation~\ref{eq:match-probability}; it applies because the event
constrains only the mode, descriptor, and sign. It also holds for
integrable signed functions by subtraction. One may include independent
numerical noises in $\Gamma$, or integrate them in $h_j$ first.

In particular the identity permits $Q_0$ to depend on the new $C_j$,
its fine gradient, marks, and all current questions. There is no
conditioning on $C_j$ in Equation~\ref{eq:match-probability}: that
partition is an output of $\Gamma$. Deleting the whole subbatch in
$B_\circ^j$ ensures that the reference sign stays fixed as $T_i$ is
redrawn. It is essential here that we do not retain the old trial's fine
partition after replacing its path data.

Apply Equation~\ref{eq:full-state-forcing} for each $i\in S$, then
average. On the resulting single actual state there is one common
$C_j$, one $K_j$, and one $\sigma$. Put
\[
 d_z=|K_{j,z}|-\sigma_{\operatorname{par}(z)}K_{j,z},\qquad
 f_m=\operatorname{avg}_{i\in S}\ind_{E_i}.
\]
Then
\begin{equation}\label{eq:match-defect-transfer}
 \operatorname{avg}_{i\in S}\E_{\rm forced}\sum_zd_z
 =\E_{\rm actual}\sum_zd_z\frac{f_m}\beta
 \le2\E\sum_zd_z+\frac2\beta\E\sum_{z\text{ marked}}|K_{j,z}|
 \le4\Delta\nu.
\end{equation}
Indeed all descriptor counts are tested, so the count for the realized
$\sigma$ is among them. On an unmarked cell $f_m\le2\beta$;
everywhere $f_m\le1$ and $0\le d_z\le2|K_{j,z}|$. This is a pointwise
empirical inequality. It does not presume any independence of a
gradient and its match indicators. Goodness and the marked-gradient
budget pay for its last two terms. In particular the unmarked defect
is not multiplied by the reciprocal match probability.

For a vertical mode, retain the hit in the empirical fraction:
$f_{m,z}^{\rm hit}=\operatorname{avg}_{i\in S}\ind_{E_i}r_i^{\rm path}(z)$.
The same calculation, now conditional on $\cZ$, gives
\begin{equation}\label{eq:match-hit-transfer}
 \operatorname{avg}_{i\in S}\E_{\rm forced}
             \left[\sum_z|K_{j,z}|r_i^{\rm path}(z)\,\middle|\,\cZ\right]
 \le2p_{\rm law}\E[|K_j|_1\mid\cZ]
       +\beta^{-1}\E\left[\sum_{z\text{ marked}}|K_{j,z}|\,
                                                 \middle|\,\cZ\right].
\end{equation}
This uses the hit-match marking rule on each unmarked cell and the
bound $f_{m,z}^{\rm hit}\le1$ elsewhere.

\subsection{Endpoint conversion and path inequalities}

In the forced experiment use $\tau=\sigma$ and $\chi=+1$.
Compared with an independent uniform descriptor, this adaptive choice
is dominated by the factor $2^J$: conditional on the remaining variables,
the latter gives that descriptor probability $2^{-J}$. The fine state
is rebuilt in both experiments by the same rule. Thus
Equation~\ref{eq:forced-endpoint-bound} implies, for either endpoint,
\begin{equation}\label{eq:adaptive-endpoint-bound}
 \operatorname{avg}_{i\in S}\E_{\rm forced}
 \frac{|h_j(B^j_{\neg i}+sD^\pm)
                      -h_{j-1}(B^j_{\neg i}+sD^\pm)|}{s}
 \le\Delta\nu.
\end{equation}
The background and the sign use none of the forced slot's fresh path
data. This domination therefore changes only the descriptor choice.

For fixed path questions, mode, descriptor, and amplitude data, the
rebuilt fine partition is fixed as $u$ varies. Along
\[
 B(u)=B^j_{\neg i}+s[(1-u)D^-+uD^+]
\]
ordinary differentiation gives
\begin{equation}\label{eq:path-derivative}
 \frac1s\frac{\dd}{\dd u}h_j(B(u))
   =\sum_zK_{j,z}(B(u))(D^+(z)-D^-(z)).
\end{equation}
Each slope is $\sigma_{\operatorname{par}(z)}w_z$ with
$0\le w_z\le2$. Its negative part is bounded, conservatively, by
$2\sum_zd_z(B(u))$: if $\sigma K\ge0$ the coordinate contributes
no negative part, whereas if $\sigma K<0$, then
$-w_z\sigma K\le2(-\sigma K)\le2d_z$.
Integrating in $u$ and using Equation~\ref{eq:full-state-forcing}
identifies the uniform-time law with the actual gradient law in
Equation~\ref{eq:match-defect-transfer}. Converting the two endpoints
with Equation~\ref{eq:adaptive-endpoint-bound} proves, for every mode,
\begin{equation}\label{eq:oriented-path-bound}
 \operatorname{avg}_{i\in S}\E_{\rm forced}
 \left(\frac{h_{j-1}(B^j_{\neg i}+sD^-)
               -h_{j-1}(B^j_{\neg i}+sD^+)}s\right)_+
 \le10\Delta\nu.
\end{equation}
Here eight units pay for fine negative variation and two for the
endpoints. The estimates are uniform in the number of slots, cells,
and descriptors.

For a vertical mode, the absolute slope in
Equation~\ref{eq:path-derivative} is at most
$2\sum_z|K_{j,z}|r(z)$. Equation~\ref{eq:match-hit-transfer}, together
with $\E[|K_j|_1\mid\cZ]\le v$, bounds its conditional expected
integral by $4p_{\rm law}v$ plus a nonnegative error whose expectation
is at most $2\Delta\nu$. The two conditional endpoint errors add at
most $2\Delta\nu$ in expectation. Conditional Jensen's inequality and
Equation~\ref{eq:coarse-width} therefore give a nonnegative function
$e_{\rm law}(\cZ)$ such that
\begin{equation}\label{eq:width-vertical}
 \E_{r\sim P_{\rm law}}|F_{\cZ}(r)|
 \le4p_{\rm law}v+e_{\rm law}(\cZ),\qquad
 \E e_{\rm law}\le4\Delta\nu.
\end{equation}
This statement also covers zero-probability laws: their expected
unmarked contribution is zero and no division by their honest
probability has been used.

For an ordered coupling $(r,r')$, write the difference of coarse widths
as the sum of the difference of their upper endpoints and the reversed
difference of their lower endpoints. The two corresponding menu entries
are $+\sigma r\to+\sigma r'$ and
$-\sigma r'\to-\sigma r$. The positive part of a sum is at most the
sum of positive parts. Applying conditional Jensen and
Equation~\ref{eq:oriented-path-bound} twice yields
\begin{equation}\label{eq:width-order}
 \E_{\cZ,r,r'}(F_{\cZ}(r)-F_{\cZ}(r'))_+\le20\Delta\nu.
\end{equation}
The same background kernel is used for both endpoints because fresh
comparison data are drawn independently after the background.

For the amplitude comparison, use the common extended kernel described
after Equation~\ref{eq:coarse-width}. At one realization abbreviate
\[
 H_a=\phi_C(B+sar),\qquad H_-=\phi_C(B-s\sigma r),\qquad
 H_+=\phi_C(B+s\sigma r).
\]
Let $a'$ be an independent fresh amplitude. For any four real numbers
of this form,
\begin{equation}\label{eq:amplitude-sandwich}
 |H_a-H_{a'}|\le H_+-H_-
   +\sum_{b\in(a,a')}
          [(H_--H_b)_++(H_b-H_+)_+].
\end{equation}
The sum counts both amplitudes, including when $a=a'$.
Indeed write
$H_a-H_{a'}=(H_a-H_+)+(H_+-H_-)+(H_--H_{a'})$ and bound the first
and last terms by their positive parts. Adding the other two nonnegative
terms gives the right side of Equation~\ref{eq:amplitude-sandwich}.
Repeating with $a,a'$ exchanged bounds the opposite signed difference
by the same expression, proving the absolute-value bound without any
ordering assumption on $H_-,H_+$.

Divide Equation~\ref{eq:amplitude-sandwich} by $s$ and average over
the background and $a'$ with $\cZ,r,a$ fixed. The absolute value of the
averaged left difference is $|g_{\cZ}(r,a)|$ by
Equation~\ref{eq:common-amplitude-kernel}; the width term averages to
$F_{\cZ}(r)$. The four positive-part errors are supplied by the two
amplitude modes, once for $a$ and once for $a'$. Integrating the iid
predicate and amplitude laws and applying
Equation~\ref{eq:oriented-path-bound} four times proves
\begin{equation}\label{eq:width-amplitude}
 \E_{\cZ,r,a}(|g_{\cZ}(r,a)|-F_{\cZ}(r))_+\le40\Delta\nu.
\end{equation}
Equations~\ref{eq:width-vertical}, \ref{eq:width-order}, and
\ref{eq:width-amplitude} have constants $4,20,40$ independent of all
later scales. In their proof the oriented-path constant was ten.

\subsection{Selecting an exterior and threshold}

Draw a threshold $\vartheta$ independently and uniformly from
$[\gamma/4,\gamma/2]$. For each exterior $(\cZ,\vartheta)$ define
the common Boolean response on actual predicates
\[
 W(r)=\ind_{\{F_{\cZ}(r)>\vartheta\}}.
\]
Equation~\ref{eq:good-stage-signal} concerns a uniform index in $I$;
Equation~\ref{eq:common-amplitude-kernel} makes its response exactly the
common $g_{\cZ}(r,a)$ under the independent iid conjunction and amplitude
laws. On $|g_{\cZ}(r,a)|\ge\gamma$ but
$F_{\cZ}(r)\le\gamma/2$, the positive part in
Equation~\ref{eq:width-amplitude} is at least $\gamma/2$. Consequently
the lost mass is at most $80\Delta\nu/\gamma$. By
Equation~\ref{eq:delta-choice},
\begin{equation}\label{eq:threshold-success}
 \E_{\cZ,\vartheta}
   [\ind_{\{v\le B_*\}}\Prob_{P_*^t}(W=1)]
 \ge b_*p\nu/2.
\end{equation}
Indeed $F>\gamma/2$ ensures success for every possible threshold, and
the event no longer depends on the auxiliary amplitude.

For every vertical law and each fixed exterior,
Equation~\ref{eq:width-vertical} and $\vartheta\ge\gamma/4$ give
\begin{equation}\label{eq:threshold-upper}
 \Prob_{\rm law}(W=1)
 \le \frac{16}{\gamma}p_{\rm law}v
                              +\frac4\gamma e_{\rm law}(\cZ).
\end{equation}
Each final term has expectation at most $16\Delta\nu/\gamma$.
For real numbers $u,u'$, the probability over this random threshold
that $\ind_{\{u>\vartheta\}}>\ind_{\{u'>\vartheta\}}$ is at most
$4(u-u')_+/\gamma$, since the relevant threshold interval has length
at most $(u-u')_+$. Thus Equation~\ref{eq:width-order} gives, for each
ordered law,
\begin{equation}\label{eq:threshold-order}
 \E_{\cZ,\vartheta}\Prob(W(r)>W(r')\mid\cZ,\vartheta)
 \le80\Delta\nu/\gamma.
\end{equation}

Let $e_{\rm tot}(\cZ,\vartheta)$ be the sum of all $V$ upper-error
terms in Equation~\ref{eq:threshold-upper} and all $Q$ conditional
violation probabilities in Equation~\ref{eq:threshold-order}. It is
nonnegative and
\begin{equation}\label{eq:total-exterior-error}
 \E e_{\rm tot}\le A_t\Delta\nu/\gamma.
\end{equation}
Discard exteriors for which $e_{\rm tot}>\eps v$. The iid law is among
the vertical laws, so Equation~\ref{eq:threshold-upper} bounds the
integrated iid success on this discarded set by
\[
 \left(\frac{16p}{\gamma\eps}+1\right)\E e_{\rm tot}
 \le b_*p\nu/4.
\]
This estimate is relative to $\nu$ and remains valid for arbitrarily
small positive cut variance. Equations~\ref{eq:threshold-success}
and~\ref{eq:total-exterior-error} leave at least $b_*p\nu/4$ integrated
success on exteriors satisfying $v\le B_*$ and
$e_{\rm tot}\le\eps v$. An exterior in this set with $v=0$ has zero
success by its iid upper bound. Moreover $\E v\le(T_0+1)\nu$ by
Equation~\ref{eq:U-mean}. Since
$b_0(T_0+1)=b_*/8$, there must be a remaining exterior with $v>0$ and
\[
 \Prob_{P_*^t}(W=1)\ge b_0vp.
\]
Otherwise its integrated success would be at most $b_*p\nu/8$.
At this same exterior, every error summand is at most $\eps v$.
Consequently every fixed type pattern and prescribed ordered law satisfy
\[
 \Prob_{\rm pat}(W=1)\le v(K_0p_{\rm pat}+\eps),\qquad
 \Prob(W(r)>W(r'))\le\eps v,
 \qquad 0<v\le B_0.
\]
This is exactly Lemma~\ref{lem:rare-event}, with a single response
depending only on the conjunction as a predicate. It produces the
valuation forbidden by Lemma~\ref{lem:event-obstruction} in a NO
instance. The contradiction proves the soundness assertion of
Theorem~\ref{thm:score-gap}.

\subsection{The interface for finite compilation}\label{subsec:score-compilation}

We record the effective structure of the construction, including its
dependence on input size. The primitive test samplers use fixed mixtures
of uniform choices from polynomial-size finite sets. They therefore have
polynomially enumerable discrete sample spaces with rational probabilities
of polynomial bit length and inverse-polynomial lower bounds on positive
seed probabilities. For fixed $T_0$, all choices in
Equation~\ref{eq:preconjunction-constants}, the conjunction length,
the path menu, $\Delta$, $M$, and then the sequential stage parameters
are constants. The upper bounds on the required integers and the
numerical precision can be arbitrarily large functions of $T_0$; none
depends on the instance size.

\begin{proposition}[Finite form structure]\label{prop:score-compilation}
The demand law and the raw comparison measures in
Theorem~\ref{thm:score-gap} have polynomially enumerable discrete seeds,
whose positive probabilities are bounded below by an inverse polynomial
in input size. For each seed, after finitely many fixed subdivisions of
numerical boxes, their conditional densities are nonnegative rational polynomials and the
score table entries are rational polynomials of bounded degree in a
fixed number of numerical variables. There are only finitely many
such numerical polynomial patterns, independent of input size.
The total raw comparison mass is finite and bounded by a constant.
For any fixed finite grid and deterministic entrywise rounding rule, every positive
rounded demand mass or comparison-pair mass has an inverse-polynomial
lower bound.
\end{proposition}

\begin{proof}
Each sampled predicate has bounded arity and a polynomially enumerable
presentation law. A conjunction, a comparison path, and each slot use
a fixed number of those samples. The number of slots, stage choices,
freeze subsets, formal descriptors, and resampling indicators is fixed.
Enumerating their products therefore takes a polynomial number of
steps. All discrete probabilities are rational, with polynomial bit
length: products involve only a fixed number of primitive probabilities
and fixed rational weights. The same observation preserves an
inverse-polynomial lower bound for each positive seed probability.
Descriptor interpretation and refining the
lists require fixed-size truth tables with positions named by the seed.
Equality of names and the values of the finite test tables can be
computed in polynomial time.

Separate the amplitude atom at zero from its uniform part. All remaining
variables are amplitudes, interpolation times, and noises in rational
intervals, together with the integration variables in the derivative
forms. Their joint densities are products of rational polynomial
densities. Absolute derivative weights become polynomial on the two
half intervals, and $-f'_{\xi_j}$ is already nonnegative on its
integration interval. No lookup-dependent descriptor occurs in any
measure. All forms are positive measures on these fixed rectangular
pieces.

Fix a seed and an assignment to its bounded list of queried positions.
Each predicate value is then zero or one, each descriptor value is a
fixed sign, and each mark is determined by the finite trial counts.
The scores are sums of rational multiples of amplitudes, times, products
of time and fresh amplitude, and noises, multiplied by these fixed
Boolean values. Hence every score entry is a polynomial of bounded
degree with rational coefficients. The possible local Boolean tables,
label-collision patterns, descriptor vectors, and mark sets all range
over fixed finite sets. The numerical coefficients $s_j,\xi_j$ and
all external form weights have already been fixed. Therefore the
polynomial patterns belong to a fixed finite family; actual bit names
only select where a table is evaluated. This also shows a uniform bound
on arity and range of every score.

Finally, each derivative integral has mass at most
$15/(8\xi_j)$ per differentiated coordinate, and there are finitely
many formal cells. Ordinary comparison measures have mass one before
their finite external weights. Summing over the fixed set of forms
gives a finite bound on total raw mass.

Fix now a finite grid and entrywise rounding rule. For one numerical pattern,
every possible rounded local table, or pair of tables, has mass equal
to the integral of its fixed polynomial density over a fixed region
of its numerical box. These regions depend only on the pattern and
the fixed rounding rule. There are finitely many patterns and rounded
tables, so the minimum of all their positive masses is a positive
constant. Renaming the queried positions or identifying collisions
does not add numerical patterns. Any positive global rounded mass has
at least one contributing discrete seed. Multiplying its
inverse-polynomial probability by this fixed positive minimum proves
the last assertion; identifying identical scores only adds masses.
This completes both the proposition and the remaining effective assertions of
Theorem~\ref{thm:score-gap}.
\end{proof}

The parameter order can now be read without circular dependencies:
$T_0$ fixes the resampling and signal constants; these fix
$e,B_*,\gamma,b_0,K_0,B_0$; the rare-event lemma fixes $t,\eps$;
then $\Delta$ fixes $M$; finally each stage fixes
$\alpha_j,m_j,n_j,s_j,\xi_j$ in that order. The path constants above
do not grow when any later scale is made smaller or any stage size
larger.
